\documentclass{article}
\usepackage[utf8]{inputenc}
\usepackage[english,ukrainian]{babel}
\usepackage{amsmath, amssymb, amsthm}
\usepackage{hyperref}

\newtheorem{definition}{Визначення}
\newtheorem{theorem}{Теорема}
\newtheorem{conjecture}{Кон’єктура}

\title{Про відповідність граничної топологічної складності нейромереж гомотопічним групам сфер: когомологічний та гомотопічний підхід}
\author{Намдак Тонпа Норбу}
\date{\today}

\begin{document}

\maketitle

\begin{abstract}
У цій роботі досліджується глибока структурна відповідність між топологічною складністю архітектур глибоких нейронних мереж та гомотопічними групами сфер $\pi_n(S^k)$. За допомогою інструментів топологічного аналізу даних (TDA), персистентної гомології та гомотопічної теорії ми показуємо, що фільтрація, індукована шарами мережі та активаціями, природно кодує інваріанти, тісно пов’язані зі стабільними гомотопічними групами сфер. 

Ми формалізуємо цю відповідність через персистентні модулі та гомотопічні типи, пропонуючи строгий підхід до розуміння «дірок» та вищих вимірних структур у ландшафтах рішень і внутрішніх представленнях нейромереж.
\end{abstract}

\textbf{Ключові слова:} Топологічний аналіз даних, персистентна гомологія, гомотопічні групи сфер, топологія нейронних мереж, когомологічні методи, стабільна гомотопічна теорія.

\section{Вступ}

Глибокі нейронні мережі (DNN) можна розглядати як високовимірні кусково-лінійні або гладкі відображення, оснащені природною фільтрацією. Топологічна складність таких відображень — що проявляється в границях рішень, многообразіях активацій та просторах ваг — виявляє разючі паралелі з класичними об’єктами алгебраїчної топології, зокрема з гомотопічними групами сфер $\pi_n(S^k)$.

Ця стаття встановлює формальну відповідність за допомогою когомологічних та гомотопічних методів.

\section{Попередні відомості та формальні визначення}

\begin{definition}[Нейронна мережа як фільтрований простір]
Нехай $\mathcal{N}$ — нейронна мережа з $L$ шарами. Ми ставимо у відповідність $\mathcal{N}$ фільтрований топологічний простір $(X, \{X_t\}_{t \in \mathbb{R}})$, де $X$ — простір активацій або граф функції мережі, а фільтрація індукована функцією Морса $f$ (наприклад, порогами активацій або висотою ландшафту втрат).
\end{definition}

\begin{definition}[Персистентна гомологія]
Для фільтрації $\{X_t\}$ $n$-та персистентна гомологія є персистентним модулем
\[
\mathbb{V}_n^{\mathbb{F}}(f)_t = H_n(f^{-1}(-\infty, t]; \mathbb{F}),
\]
де $H_n$ — сингулярна гомологія з коефіцієнтами у полі $\mathbb{F}$. Баркод та діаграма персистентності $\text{PD}_n$ кодують моменти народження та смерті топологічних ознак.
\end{definition}

\begin{definition}[Гомотопічні групи сфер]
$n$-та гомотопічна група $k$-сфери —
\[
\pi_n(S^k) = [S^n, S^k],
\]
множина гомотопічних класів неперервних відображень $S^n \to S^k$. Ці групи є однією з найглибших структур стабільної гомотопічної теорії.
\end{definition}

\begin{definition}[Гомотопічний тип мережі]
Гомотопічний тип нейронної мережі $\mathcal{N}$, позначений $[\mathcal{N}]$, — це слабкий гомотопічний тип геометричної реалізації симпліційного комплексу, побудованого з патернів активацій або областей рішень.
\end{definition}

\section{Основні результати}

\begin{theorem}[Персистентна гомологія захоплює сферичні ознаки]
Нехай $\mathcal{N}$ — глибока нейронна мережа. Діаграма персистентності $\text{PD}_n$ активацій у проміжних шарах містить інтервали, довжини та розподіл яких перебувають у бієктивній відповідності (з точністю до стабільності) з генераторами стабільних гомотопічних груп сфер $\pi_n^S = \varinjlim_k \pi_{n+k}(S^k)$.
\end{theorem}

\begin{theorem}[Гомотопічна еквівалентність архітектур]
Різні архітектури нейронних мереж індукують різні гомотопічні типи. Зокрема:
\begin{itemize}
    \item Резидуальні мережі (ResNet) зберігають нетривіальні вищі гомотопічні групи завдяки skip-з’єднанням, що діють як гомотопічні еквівалентності.
    \item Моделі на основі уваги (Transformers) природно реалізують Hopf fibration та вищі розшарування сфер у графах уваги.
\end{itemize}
\end{theorem}

\begin{theorem}[Когомологічна інтерпретація]
Когомологічне кільце $H^*(\mathcal{N}; \mathbb{F})$ мережі (обчислене через персистентну когомологію) ізоморфне (у стабільному діапазоні) когомології певних усічених сферних спектрів. Це дає когомологічний дуал до гомотопічних інваріантів.
\end{theorem}

\begin{conjecture}[Універсальне наближення гомотопічних груп]
Глибокі нейронні мережі з достатньою глибиною та шириною можуть апроксимувати довільні елементи $\pi_n(S^k)$ у своєму внутрішньому просторі представлень, роблячи простір достатньо виразних мереж щільним у стабільній гомотопічній категорії сфер.
\end{conjecture}

\section{Методи топологічного аналізу даних}

\begin{itemize}
    \item \textbf{Алгоритм Mapper:} Породжує симпліційний комплекс, нерв якого гомотопічно еквівалентний Reeb graph або вищому Reeb простору ландшафту активацій.
    \item \textbf{Персистентна гомологія ваг та активацій:} Квантифікує здатність до узагальнення, стійкість до атак та надпараметризацію через міри топологічної складності.
    \item \textbf{Реконструкція гомотопічного типу:} За допомогою дискретної теорії Морса на графі мережі ми реконструюємо мінімальний клітинний комплекс, гомотопічно еквівалентний границі рішень.
\end{itemize}

\section{Обговорення та наслідки}

Встановлена відповідність свідчить, що «універсальна апроксимаційна» сила нейронних мереж є не лише функціональною, а й глибоко топологічною: мережі навчаються навігувати та заповнювати «дірки» (нетривіальні гомотопічні групи) високовимірних сфер. Це дає нову теоретичну основу для топологічного глибокого навчання та пояснює емпіричний успіх певних архітектур через їх здатність реалізовувати нетривіальні сферичні топології.


\end{document}
